% Bagian: first-order-logic
% Bab: proof-systems
% Subbagian: introduction

\documentclass[../../../include/open-logic-section]{subfiles}

\begin{document}

\iftag{FOL}
      {\olfileid[id]{fol}{prf}{int}}
      {\olfileid[id]{pl}{prf}{int}}

\olsection{Pendahuluan}

Logika pada umumnya memiliki semantik sekaligus sistem !!a{derivation}.
Semantik berkenaan dengan konsep-konsep seperti kebenaran,
keterpenuhan, validitas, dan perikutan. Tujuan sistem !!{derivation}
adalah menyediakan metode yang sepenuhnya sintaktis untuk menetapkan
perikutan dan validitas. Sistem itu sepenuhnya sintaktis dalam arti
bahwa suatu !!{derivation} di dalam sistem demikian merupakan objek
sintaktis berhingga, biasanya berupa barisan (atau susunan berhingga
lainnya) dari !!{sentence}s atau !!{formula}s. Sistem !!{derivation}
yang baik memiliki sifat bahwa setiap barisan atau susunan tertentu
dari !!{sentence}s atau !!{formula}s dapat diperiksa secara mekanis
apakah ``benar.''

Sistem !!{derivation} yang paling sederhana (dan secara historis
pertama) bagi logika orde pertama bersifat \emph{aksiomatik}. Sebuah
barisan !!{formula}s dianggap sebagai !!a{derivation} dalam sistem
demikian apabila setiap !!{formula} di dalamnya merupakan salah satu
dari sekumpulan tetap ``aksioma'' atau mengikuti dari !!{formula}s
yang mendahuluinya dalam barisan melalui salah satu dari sejumlah
tetap ``aturan inferensi''---dan dapat diperiksa secara mekanis apakah
!!a{formula} merupakan aksioma serta apakah formula itu mengikuti
secara tepat dari !!{formula}s lain melalui salah satu aturan
inferensi. Sistem !!{derivation} aksiomatik mudah dideskripsikan---dan
juga mudah ditangani secara metateoretis---tetapi !!{derivation}s di
dalamnya sulit dibaca dan dipahami, serta juga sulit dihasilkan.

Sistem !!{derivation} lain telah dikembangkan dengan tujuan agar
!!{derivation}s lebih mudah disusun atau !!{derivation}s lebih mudah
dipahami setelah
selesai. Contohnya ialah deduksi alami, pohon kebenaran, yang juga
dikenal sebagai pembuktian tableau, dan kalkulus sekuen. Sejumlah
sistem !!{derivation} dirancang khusus dengan mempertimbangkan
mekanisasi; misalnya, metode resolusi mudah diimplementasikan dalam
perangkat lunak (tetapi !!{derivation}s di dalamnya pada dasarnya mustahil
dipahami). Sebagian besar sistem !!{derivation} lain ini
merepresentasikan !!{derivation}s sebagai pohon !!{formula}s, bukan
sebagai barisan. Dengan demikian, lebih mudah dilihat bagian mana dari
!!a{derivation} yang bergantung pada bagian lainnya.

Jadi, untuk suatu logika, seperti logika orde pertama, berbagai sistem
!!{derivation} akan memberikan eksplikasi yang berbeda mengenai apa
artinya !!a{sentence} merupakan \emph{teorema} dan apa artinya
!!a{sentence} !!{derivable} dari beberapa kalimat lain. Apa pun
caranya (melalui !!{derivation}s aksiomatik, deduksi alami,
!!{derivation}s sekuen, pohon kebenaran, atau refutasi resolusi), kita
menginginkan agar relasi-relasi ini bersesuaian dengan gagasan
semantis validitas dan perikutan. Kita tulis $\Proves !A$ untuk
``$!A$~merupakan teorema'' dan ``$\Gamma \Proves !A$'' untuk
``$!A$~!!{derivable} dari~$\Gamma$.'' Bagaimanapun $\Proves$~
didefinisikan, kita menginginkannya bersesuaian dengan $\Entails$,
yaitu:
\begin{enumerate}
\item $\Proves !A$ jika dan hanya jika $\Entails !A$
\item $\Gamma \Proves !A$ jika dan hanya jika $\Gamma \Entails !A$
\end{enumerate}
Arah ``hanya jika'' di atas disebut \emph{kesahihan}.
Suatu sistem !!{derivation} bersifat sahih jika !!{derivability} menjamin
perikutan (atau validitas). Setiap sistem !!{derivation} yang layak
harus bersifat sahih; sistem !!{derivation} yang tidak sahih sama
sekali tidak berguna. Lagi pula, seluruh tujuan !!a{derivation} ialah
memberikan jaminan sintaktis atas validitas atau perikutan. Kita akan
membuktikan kesahihan bagi sistem-sistem !!{derivation} yang kita
sajikan.

Arah kebalikan, yakni arah ``jika,'' juga penting: arah ini disebut
\emph{kelengkapan}. Sistem !!{derivation} yang lengkap cukup kuat
untuk menunjukkan bahwa $!A$~merupakan teorema setiap kali $!A$~valid,
dan bahwa $\Gamma \Proves !A$ setiap kali $\Gamma \Entails !A$.
Kelengkapan lebih sulit dibuktikan, dan beberapa logika tidak memiliki
sistem !!{derivation} yang lengkap. Logika orde pertama memilikinya.
Kurt G\"odel adalah orang pertama yang membuktikan kelengkapan bagi
sistem !!a{derivation} logika orde pertama dalam disertasinya pada
tahun 1929.

Konsep lain yang berkaitan dengan sistem !!{derivation} ialah
\emph{konsistensi}. Himpunan !!{sentence}s disebut tidak konsisten jika
apa pun dapat diturunkan darinya, dan disebut konsisten jika tidak
demikian. Ketidakkonsistenan merupakan padanan sintaktis dari
ketakterpenuhan: seperti himpunan yang tidak dapat dipenuhi, himpunan
!!{sentence}s yang tidak konsisten tidak dapat menjadi teori yang
baik; himpunan tersebut cacat secara mendasar. Himpunan !!{sentence}s
yang konsisten belum tentu benar atau berguna, tetapi setidaknya
memenuhi ambang minimal kegunaan logis tersebut. Untuk sistem
!!{derivation} yang berbeda, definisi khusus konsistensi himpunan
!!{sentence}s mungkin berbeda, tetapi seperti~$\Proves$, kita
menginginkan agar konsistensi bertepatan dengan padanan semantisnya,
yaitu keterpenuhan. Kita menginginkan agar selalu berlaku bahwa
$\Gamma$ konsisten jika dan hanya jika ia dapat dipenuhi. Di sini,
arah ``hanya jika'' sama artinya dengan kelengkapan (konsistensi
menjamin keterpenuhan), sedangkan arah ``jika'' sama artinya dengan
kesahihan (keterpenuhan menjamin konsistensi). Bahkan, untuk logika
orde pertama klasik, kedua versi kesahihan dan kelengkapan itu
ekuivalen.

\end{document}
